\code{torch.compile} lowers a PyTorch model through TorchInductor, and a large part of what Inductor does is graph
rewriting: a few hundred rules each look for a pattern (a manual attention computation, a convolution followed by
batch normalisation, a matrix multiplication of a particular shape) and replace it with something the rule's author
expected to be faster. Rules fire whenever they match. Nothing measures whether the result is faster on the target,
and when a compiled model is slower than expected there is no way, for users or for the compiler's developers, to
find out which rule did it~\cite{rfc}. The only published per-optimization view of Inductor is an eight-knob
ablation on one GPU reported as a geometric mean~\cite{pt2}.

This paper makes the question answerable and reports what the answer looks like on six machine/precision targets.
Our thesis is that a graph rewrite fixes the \emph{shape} of the generated program, and whether that shape is faster
depends on the relative cost of library calls, generated kernels and memory traffic on the target, which no rule
written once can know. Three measured facts support it: half of the most impactful rules are confirmed to help on
one machine and hurt on another (\S\ref{sec:signs}); when they flip, the compiler has made the identical structural
change on both machines (Table~\ref{tab:mechanism}); and acting on the measurement pays, either through a learned
guard per rule where harm is frequent and predictable or through a per-machine default where it is not
(\S\ref{sec:guards}).

\noindent\textbf{Contributions.}
\begin{enumerate}[leftmargin=*]
\item \tool: per-rule switches for every pattern rule, configuration flag and opt-in pass of TorchInductor 2.14
  (269/396/398 switches on the Apple, x86 and CUDA builds), installed from outside the compiler by wrapping the hook it
  already calls before applying a rule, with fire and suppression counts (\S\ref{sec:switches}).
\item A timing oracle whose false-alarm rate was measured on 480 runs of unchanged configurations and reduced from
  9\% to 1.5\% by two time-separated confirmations, and delta-debugging attribution over switch sets that recovers
  planted culprits in 5--16 oracle calls, detects pairs, and declines to blame when nothing is slower
  (\S\ref{sec:attr}).
\item A study of 24 models on an Apple M4 Pro, an Intel Xeon, and NVIDIA A40, L4 and RTX~4090 GPUs in fp32 and bf16:
  the rule$\times$machine sign matrix, its mechanism, and the effect of precision on which rules fire at all
  (\S\ref{sec:signs}).
\item Two ways to act on blame, both evaluated on held-out models against a per-model oracle: learned per-rule guards,
  and per-machine defaults; the latter yields a concrete upstream finding about \code{freezing} backed by 163 real
  networks per GPU, and a Dynamo bug (\S\ref{sec:guards}).
\end{enumerate}
